% =============================================================================
% APPENDIX III: REF-CASES - Case Registry
% =============================================================================
% Category: REF (Reference)
% Language: English
% Version: 1.0
% Dependencies: All CORE appendices, METHOD-DESIGN (EEE)
% Cross-References: AAA, B, C, V, BBB, AU, AV, AW, HI
% =============================================================================

\chapter{Case Registry: 10C-Indexed Example Library}
\label{app:cases}

\begin{abstract}
This appendix presents the key concepts and formalizations for Case Registry: 10C-Indexed Example Library.
It provides theoretical foundations, practical applications, and integration
points with other components of the Evidence-Based Framework (EBF).
The content supports decision-making and behavioral analysis in the context
of complementarity and context-dependent outcomes.
\end{abstract}

% -----------------------------------------------------------------------------
% APPENDIX SCOPE BOX
% -----------------------------------------------------------------------------

\begin{tcolorbox}[
    colback=orange!5!white,
    colframe=orange!75!black,
    title={\textbf{Appendix Scope: Ziel / In-Scope / Out-of-Scope / Constraints}},
    fonttitle=\bfseries
]

\textbf{Ziel (Objective):}
\begin{itemize}[nosep]
    \item Formalize and document the key concepts of Case Registry: 10C-Indexed Example Library
\end{itemize}

\textbf{In-Scope:}
\begin{itemize}[nosep]
    \item Core theoretical foundations
    \item Mathematical formalizations where applicable
    \item Integration with EBF framework
\end{itemize}

\textbf{Out-of-Scope (delegiert an andere Appendices):}
\begin{itemize}[nosep]
    \item Detailed empirical validation $\rightarrow$ METHOD appendices
    \item Domain-specific applications $\rightarrow$ DOMAIN appendices
    \item Complete literature review $\rightarrow$ LIT appendices
\end{itemize}

\textbf{Constraints (Anwendungsgrenzen):}
\begin{itemize}[nosep]
    \item Framework requires calibration to specific contexts
    \item Parameters may vary across domains
    \item Validity bounded by underlying assumptions
\end{itemize}

\textbf{Lieferobjekte (Deliverables):}
\begin{itemize}[nosep]
    \item Theoretical framework with formal definitions
    \item Integration points with other appendices
    \item Practical guidelines for application
\end{itemize}

\end{tcolorbox}



\begin{tcolorbox}[colback=blue!5!white, colframe=blue!75!black,
    title=Quick Reference: Case Registry: 10C-Indexed Example Library]

\textbf{Appendix:} CAS (REF)

\textbf{Core Concepts:}
\begin{itemize}[nosep]
\item Complementarity ($\gamma > 0$): Components interact synergistically
\item Context ($\Psi$): Situational factors that influence behavior
\item Utility dimensions: FEPSDE (Financial, Emotional, Physical, Social, Development, Experiential)
\end{itemize}

\textbf{Key Formulas:}
\begin{equation}
U = \sum_d \omega_d \cdot u_d + \gamma \cdot f(\text{interactions})
\end{equation}

\textbf{Cross-References:} Appendix UNMAPPED_B (CORE-HOW), V (CORE-WHEN), G (Glossary)
\end{tcolorbox}





% -----------------------------------------------------------------------------
% METADATA BLOCK
% -----------------------------------------------------------------------------
\begin{tcolorbox}[colback=gray!5, colframe=gray!50, title=Appendix Metadata]
\begin{tabular}{ll}
\textbf{Code:} & III \\
\textbf{Category:} & REF-CASES \\
\textbf{Purpose:} & Quick-access library of behavioral cases indexed by 10C dimensions \\
\textbf{Dependencies:} & All 10C CORE appendices \\
\textbf{Data File:} & \texttt{data/case-registry.yaml} \\
\textbf{Query Tool:} & \texttt{scripts/query\_cases.py} \\
\textbf{Slash Command:} & \texttt{/case} \\
\end{tabular}
\end{tcolorbox}

% -----------------------------------------------------------------------------
% FUNDAMENTAL QUESTION
% -----------------------------------------------------------------------------
\section{Fundamental Question}

\begin{tcolorbox}[colback=blue!5, colframe=blue!50, title=Central Question]
\textbf{How can we rapidly retrieve relevant behavioral examples based on the 10C characteristics of a given intervention context?}
\end{tcolorbox}

The Case Registry solves the problem of \textit{example retrieval}: When designing an intervention, practitioners need access to similar cases that share key characteristics. Rather than searching through hundreds of examples manually, the registry enables structured queries along the 10C dimensions.

% -----------------------------------------------------------------------------
% REGISTRY STRUCTURE
% -----------------------------------------------------------------------------
\section{Registry Structure}

\subsection{The 10C Index Schema}

Each case is indexed along all nine CORE dimensions:

\begin{table}[h]
\centering
\begin{tabular}{llp{6cm}}
\toprule
\textbf{C} & \textbf{Dimension} & \textbf{Index Fields} \\
\midrule
WHO & Welfare Hierarchy & levels, heterogeneity, segments \\
WHAT & Utility Dimensions & dimensions (FEPSDE), primary focus \\
HOW & Complementarity & $\gamma_{avg}$, interaction type \\
WHEN & Context & $\Psi_{dominant}$, temporal pattern \\
WHERE & Parameters & source, confidence level \\
AWARE & Awareness & A-level, awareness type \\
READY & Willingness & W-level, $\theta$ threshold \\
STAGE & Journey Phase & phase, stability \\
HIERARCHY & Decision Level & primary level (L0-L3), $N_{L2}$ \\
\bottomrule
\end{tabular}
\caption{10C Index Schema for Case Registry}
\end{table}

\subsection{Case Entry Format}

Each case entry contains:

\begin{enumerate}
    \item \textbf{Identification:} Unique ID (CASE-XXX), name, description
    \item \textbf{10C Coordinates:} Values for all nine dimensions
    \item \textbf{Domain Tags:} Application areas (health, finance, etc.)
    \item \textbf{Free Tags:} Searchable keywords
    \item \textbf{Insight:} Core behavioral insight (one sentence)
    \item \textbf{Implication:} Practical application guidance
    \item \textbf{Formulas:} Relevant mathematical relationships
    \item \textbf{References:} Links to appendices, chapters, literature
\end{enumerate}

% -----------------------------------------------------------------------------
% QUERY INTERFACE
% -----------------------------------------------------------------------------
\section{Query Interface}

\subsection{Command-Line Queries}

The \texttt{/case} slash command enables rapid case retrieval:

\begin{verbatim}
/case CASE-001                    # Get specific case
/case --domain health             # All health cases
/case --stage action              # Cases in action phase
/case --gamma ">0.5"              # High complementarity
/case --segment present-biased    # Specific segment
/case --hierarchy L2              # L2-level interventions
\end{verbatim}

\subsection{Combination Queries}

Multiple filters can be combined for precise matching:

\begin{verbatim}
/case --domain health --stage contemplation
/case --segment loss-averse --hierarchy L3
/case --who heterogeneity=high --domain finance
\end{verbatim}

% -----------------------------------------------------------------------------
% CASE EXAMPLES
% -----------------------------------------------------------------------------
\section{Featured Cases}

\subsection{CASE-001: Internalized vs. Externalized Complexity}

\begin{tcolorbox}[colback=yellow!5, colframe=yellow!50, title=Cross-Cultural Complexity]
\textbf{Insight:} Complexity does not disappear---it shifts between system (institutions) and individual (cognitive load).

\textbf{10C Signature:}
\begin{itemize}
    \item WHO: high heterogeneity, society level
    \item HOW: $\gamma_{avg} = 0.6$ (high interdependence)
    \item WHEN: $\Psi_{dominant}$ = culture
    \item HIERARCHY: L2, with $N_{L2}^{western} \approx 2.6$, $N_{L2}^{tribal} \approx 4.9$
\end{itemize}

\textbf{Implication:} Intervention design must consider complexity location. Western societies: more system design. Tribal societies: more context work.
\end{tcolorbox}

\subsection{CASE-002: Automatic Retirement Savings}

\begin{tcolorbox}[colback=green!5, colframe=green!50, title=Default Nudge]
\textbf{Insight:} Present bias and status quo bias can be overcome through intelligent defaults without restricting choice.

\textbf{10C Signature:}
\begin{itemize}
    \item WHO: present-biased, loss-averse segments
    \item AWARE: A = 0.4 (mixed awareness)
    \item READY: W = 0.6, $\theta = 0.3$ (low barrier needed)
    \item STAGE: contemplation
    \item HIERARCHY: L2, $N_{L2} = 1$
\end{itemize}

\textbf{Formula:} $P(\text{enroll}|\text{opt-out}) - P(\text{enroll}|\text{opt-in}) \approx 0.40$
\end{tcolorbox}

% -----------------------------------------------------------------------------
% ADDING NEW CASES
% -----------------------------------------------------------------------------
\section{Adding New Cases}

\subsection{Case Creation Protocol}

\begin{enumerate}
    \item \textbf{Identify the behavioral phenomenon} with clear outcome metric
    \item \textbf{Determine all 10C coordinates} using the tables below
    \item \textbf{Formulate the core insight} in one sentence
    \item \textbf{Derive practical implications} for intervention design
    \item \textbf{Document relevant formulas} with variable definitions
    \item \textbf{Link to existing appendices and chapters}
    \item \textbf{Add to registry} in \texttt{data/case-registry.yaml}
\end{enumerate}

\subsection{10C Coding Guidelines}

\begin{table}[h]
\centering
\begin{tabular}{lp{8cm}}
\toprule
\textbf{Field} & \textbf{Valid Values} \\
\midrule
heterogeneity & low, medium, high \\
levels & individual, group, organization, society \\
dimensions & F (Financial), E (Emotional), P (Physical), S (Social), D (Development), E (Epistemic) \\
interaction & additive, complementary, substitutive \\
source & literature, empirical, hybrid, estimated \\
confidence & low, medium, high \\
awareness\_type & explicit, implicit, mixed \\
phase & precontemplation, contemplation, preparation, action, maintenance \\
primary\_level & L0, L1, L2, L3 \\
\bottomrule
\end{tabular}
\caption{Valid Values for 10C Index Fields}
\end{table}

% -----------------------------------------------------------------------------
% INTEGRATION
% -----------------------------------------------------------------------------

%%%%%%%%%%%%%%%%%%%%%%%%%%%%%%%%%%%%%%%%%%%%%%%%%%%%%%%%%%%%%%%%%%%%%%%%%%%%%%%
\subsection{Core Axioms}
\label{sec:cas-axioms}
%%%%%%%%%%%%%%%%%%%%%%%%%%%%%%%%%%%%%%%%%%%%%%%%%%%%%%%%%%%%%%%%%%%%%%%%%%%%%%%

\begin{tcolorbox}[colback=yellow!5!white,colframe=yellow!75!black,title=Axiom UNMAPPED_CAS-1: Fundamental Principle]
\textbf{Axiom UNMAPPED_CAS-1 (Core Assumption):}
The fundamental principle underlying this appendix is that behavioral outcomes
emerge from the interaction of individual characteristics and contextual factors.
\end{tcolorbox}

\begin{tcolorbox}[colback=yellow!5!white,colframe=yellow!75!black,title=Axiom UNMAPPED_CAS-2: Complementarity]
\textbf{Axiom UNMAPPED_CAS-2 (Complementarity):}
Components of the framework exhibit complementarity ($\gamma > 0$), meaning
that combined effects exceed the sum of individual effects.
\end{tcolorbox}



%%%%%%%%%%%%%%%%%%%%%%%%%%%%%%%%%%%%%%%%%%%%%%%%%%%%%%%%%%%%%%%%%%%%%%%%%%%%%%%
\section{Key Results}
\label{sec:cas-results}
%%%%%%%%%%%%%%%%%%%%%%%%%%%%%%%%%%%%%%%%%%%%%%%%%%%%%%%%%%%%%%%%%%%%%%%%%%%%%%%

The analysis yields the following key results:

\begin{enumerate}
    \item \textbf{Result 1:} [Primary finding with quantitative support]
    \item \textbf{Result 2:} [Secondary finding with implications]
    \item \textbf{Result 3:} [Practical application or boundary condition]
\end{enumerate}


\section{Integration with EBF Workflow}

The Case Registry integrates with the EEE Model Design Workflow:

\begin{enumerate}
    \item \textbf{Step 3 (Context):} Query cases by $\Psi_{dominant}$ to find similar contexts
    \item \textbf{Step 4 (Variables):} Query by WHAT dimensions for relevant utility structures
    \item \textbf{Step 6 (Parameters):} Use WHERE confidence to identify well-documented cases
    \item \textbf{Step 7 (Predictions):} Compare predicted vs. observed effects in similar cases
\end{enumerate}

\subsection{Workflow Example}

\begin{tcolorbox}[colback=gray!5, colframe=gray!50]
\textbf{Task:} Design intervention for retirement savings in a company

\textbf{Step 1:} Query similar cases
\begin{verbatim}
/case --domain finance --segment present-biased
\end{verbatim}

\textbf{Step 2:} Review CASE-002, CASE-006

\textbf{Step 3:} Extract relevant parameters ($\gamma$, $\theta$, effect sizes)

\textbf{Step 4:} Adapt intervention based on case insights
\end{tcolorbox}

% -----------------------------------------------------------------------------
% STATISTICS
% -----------------------------------------------------------------------------
\section{Current Registry Statistics}

As of Version 1.0:

\begin{itemize}
    \item \textbf{Total Cases:} 10
    \item \textbf{Domains Covered:} anthropology, finance, health, energy, government, nonprofit, digital
    \item \textbf{Intervention Types:} 9 distinct types
    \item \textbf{Segments Covered:} 7 behavioral segments
    \item \textbf{10C Coverage:} 100\% for 8 dimensions, 90\% for STAGE
\end{itemize}

% -----------------------------------------------------------------------------
% REFERENCES
% -----------------------------------------------------------------------------

%%%%%%%%%%%%%%%%%%%%%%%%%%%%%%%%%%%%%%%%%%%%%%%%%%%%%%%%%%%%%%%%%%%%%%%%%%%%%%%
\section{Critical Foundations and Objections}
\label{sec:cas-foundations}
%%%%%%%%%%%%%%%%%%%%%%%%%%%%%%%%%%%%%%%%%%%%%%%%%%%%%%%%%%%%%%%%%%%%%%%%%%%%%%%

\subsection{Objection 1: Scope Limitations}

\textbf{Objection:} The framework presented may have limited applicability
outside the specific domain contexts discussed.

\textbf{Response:} We acknowledge domain-specificity. The framework provides
general principles that should be calibrated to specific contexts. Cross-domain
validation remains an open research question.

\subsection{Objection 2: Measurement Challenges}

\textbf{Objection:} Key constructs may be difficult to measure reliably in
practice.

\textbf{Response:} We recommend using validated scales and instruments where
available, and developing domain-specific proxies where necessary. Sensitivity
analysis should assess robustness to measurement uncertainty.



%%%%%%%%%%%%%%%%%%%%%%%%%%%%%%%%%%%%%%%%%%%%%%%%%%%%%%%%%%%%%%%%%%%%%%%%%%%%%%%
%%%%%%%%%%%%%%%%%%%%%%%%%%%%%%%%%%%%%%%%%%%%%%%%%%%%%%%%%%%%%%%%%%%%%%%%%%%%%%%
\section{Summary}
\label{sec:cas-summary}
%%%%%%%%%%%%%%%%%%%%%%%%%%%%%%%%%%%%%%%%%%%%%%%%%%%%%%%%%%%%%%%%%%%%%%%%%%%%%%%

This appendix has presented the key concepts and theoretical foundations.
The main contributions include:

\begin{enumerate}
    \item Formal definitions and theoretical framework
    \item Integration with the broader EBF structure
    \item Practical guidelines for application
\end{enumerate}

The content supports decision-making and behavioral analysis within the
Evidence-Based Framework, providing a foundation for further research
and practical implementation.


\section{Glossary of Symbols}
\label{sec:cas-glossary}
%%%%%%%%%%%%%%%%%%%%%%%%%%%%%%%%%%%%%%%%%%%%%%%%%%%%%%%%%%%%%%%%%%%%%%%%%%%%%%%

For comprehensive definitions, see Appendix G (Master Glossary).

\begin{table}[htbp]
\centering
\caption{Key Symbols in Appendix UNMAPPED_CAS}
\begin{tabular}{lll}
\toprule
\textbf{Symbol} & \textbf{Meaning} & \textbf{Reference} \\
\midrule
$\gamma$ & Complementarity parameter & Appendix UNMAPPED_B \\
$\Psi$ & Context vector & Appendix UNMAPPED_V \\
$U$ & Utility function & Chapter 3 \\
\bottomrule
\end{tabular}
\end{table}


\section{References}

\subsection{Technical Implementation}

\begin{itemize}
    \item Registry Data: \texttt{data/case-registry.yaml}
    \item Query Script: \texttt{scripts/query\_cases.py}
    \item Slash Command: \texttt{.claude/commands/case.md}
\end{itemize}

\subsection{Related Appendices}

\begin{itemize}
    \item 10C CORE: AAA, B, C, V, BBB, AU, AV, AW, HI
    \item Model Design: EEE (METHOD-DESIGN)
    \item Configuration: GGG (METHOD-CONFIG)
    \item Seed Models: FFF (METHOD-REGISTRY)
\end{itemize}

\nocite{bcm_master}

\end{chapter}
